\section[Complete local tree construction]{A complete local tree construction from the top-tree interface}
\label{sec:op3-local-top-trees}

\paragraph{Status and scope.}
This section discharges the hierarchy condition in
\cref{sec:op3-ordered-path-clusters} for arbitrary finite unweighted trees.
It is \statusproved{} as a proof draft awaiting independent review. The
published balancing algorithm is a \statussource{} import; the application
callbacks and exact audits are implemented here. The exhaustive audit driver
is not a fast online balancing implementation. No conclusion for arbitrary
cyclic graphs follows from this tree result.
Use the exact-real word model and the local degree/adjacency access model of
\cref{sec:scope}. The physical seed is $v$,
$\boldsymbol M=\boldsymbol D-\gamma\boldsymbol A$, and
$\boldsymbol b=\boldsymbol e_v-\lambda\boldsymbol d$, with
$\lambda=\epsappr/2$. Every checkpoint exposes $v$ as the single external
boundary of the current active tree. Depths and parent edges refer to that
permanently rooted mathematical tree, not to the changing cluster hierarchy.

\begin{lemma}[Rooted ports and unavailable summaries; \statusproved]
\label{lem:op3-source-interior-adapter}
During a leaf insertion and subsequent exposure of $v$, a cluster can store
its port set, a Boolean indicating whether it contains $v$, and either its
exact summary from \cref{sec:op3-ordered-path-clusters} or an unavailable
marker. Mark it unavailable exactly when $v$ is in its interior.
A valid cluster has one or two ports, its rootward vertex is a port, and its
two ports, if present, are comparable in the tree rooted at $v$.
A valid parent never needs an unavailable child. After exposing $v$, every
cluster is valid. Marking, splitting and propagating these flags cost
constant work per structural operation.
\end{lemma}
\begin{proof}
The maintained forest has one component containing edges, namely the active
tree containing $v$, and at most one newly allocated isolated vertex.
A zero-port edge cluster must be the whole component, hence contains $v$
in its interior and is unavailable. Consider any other cluster $C$ whose
rootward vertex is $r$. If $v\notin C$, the edge from $r$ toward $v$ leaves
$C$, making $r$ a boundary. If $v\in C$ and $C$ is valid, then $r=v$ is a
boundary by validity. Thus the rootward vertex is always a port. With at
most two ports, they are comparable: otherwise their least common ancestor,
or the exit from it toward $v$, would give a third boundary.

A vertex interior to a child cannot be a boundary of its parent: every edge
incident to it already belongs to the child, and it is not externally
exposed. In particular, an unavailable child forces its parent to be
unavailable. Conversely, a valid parent has only valid children and may use
their algebraic summaries. These observations apply between updates.
The source discipline in Section 2.1 of \cite{alstrup2005toptrees} first
splits every cluster whose edges or boundary set will change. Thus an old
interior vertex cannot silently become a new boundary in a retained cluster.
After exposure, every cluster containing $v$ has it as a boundary, so every
unavailable cluster has been split or its base edge recreated.
Containment is the Boolean OR of the children's flags; no membership scan
is needed. An unavailable node retains its structural child pointers so a
later split needs no algebraic reconstruction.
\end{proof}

\begin{lemma}[Complete callback closure; \statusproved]
\label{lem:op3-top-tree-callback-closure}
With the convention of \cref{lem:op3-source-interior-adapter}, every base
creation and legal top-tree join is implemented by a constant number of the
edge, compress, rake and forget operations of
\cref{sec:op3-ordered-path-clusters}. It costs
$O(\log^3(2+N))$ exact word work and allocates
$O(\log^2(2+N))$ hull nodes, where $N$ bounds the number of owned rows.
No source load is eliminated in a valid summary.
\end{lemma}
\begin{proof}
At a base edge, sort its endpoints by stored depth. With two ports, use the
edge summary; with one port, retain its rootward endpoint and forget its
nonseed farther endpoint. A source-interior base edge is unavailable and
requires no algebra.

In a join, the two child clusters share exactly one vertex, which is a port
of each child. Every other child port remains a parent port: an edge leaving
the child there cannot belong to the other child, since their only common
vertex is the shared port. There are therefore exactly the following cases.
\begin{enumerate}
\item Two two-port children produce a two-port parent and eliminate their
shared nonseed middle port. Their depth ordering identifies left and right,
so \cref{lem:op3-adjacent-chain-separation} gives a separated compression.
\item One one-port and one two-port child produce either a two-port parent,
by raking, or a one-port parent, by raking and forgetting the shared port.
In the latter case the shared port is farther from $v$; otherwise the
parent's rootward vertex would be interior, contrary to
\cref{lem:op3-source-interior-adapter}.
\item Two one-port children share their sole port. If it remains a parent
port, rake their scalar states and take the minimum threshold. Otherwise
the parent has no ports and is unavailable.
\end{enumerate}
An unavailable parent is constructed as a structural record, independently
of its children's algebraic status. For every valid parent all children
are valid by \cref{lem:op3-source-interior-adapter}. The list exhausts the
source's five join shapes. At a valid elimination the eliminated vertex is
not $v$, and the eliminated interior is source-free. Its negative load
invariant and positive Schur pivots therefore hold. The primitive bounds in
\cref{lem:op3-persistent-projective-chain} finish the accounting.
\end{proof}

\begin{lemma}[One edge payload and online allocation; \statusproved]
\label{lem:op3-top-tree-payload}
A changed owned boundary row on one base edge can be maintained in
$O(\log^4(2+n))$ word work and $O(\log^3(2+n))$ allocated hull nodes in a
height-bounded top tree on $n$ active vertices. The local application needs
no preprocessing indexed by the ambient graph size.
\end{lemma}
\begin{proof}
Use the parent/child and edge-leaf pointers in the source representation,
and walk from the changed leaf to the current hierarchy root. Recompute the
base summary and every ancestor from its current children, preserving its
existing port set and structural join shape. Even unavailable ancestors
receive fresh child pointers. The walk has logarithmic length by Theorem 2.1 (journal p.\ 247)
of \cite{alstrup2005toptrees}; the callback bound proves the claim. This
operation changes application data only, so it does not rebalance the
structural hierarchy. Later joins, including joins that undo a previous
exposure, always derive fresh summaries from current children. They do not
reactivate stale numerical snapshots from before the payload change.

The underlying forest contains only admitted vertices and one newly created
isolated vertex. Source vertex and edge identifiers can therefore be local
indices allocated on demand. Store a comparison dictionary from discovered
original labels to local records, and keep each active incidence list as
local edge records. Resizable arrays can instead be allocated by doubling,
charging each copy to the growing number of exposed words. The source's
internal treatment of large degrees concerns these active structural
incidences; it neither accesses inactive adjacency nor adds vertices to the
PageRank matrix. Initialization of a singleton's empty top tree needs
constant state. The dictionary, cache, allocation and occasional copies all
cost at most a logarithmic factor per discovered word.

If a concrete source implementation requests a fixed vertex capacity,
start with constant capacity and double it when exhausted. Rebuild only
the currently active structural tree, using its cached edges and current
home rows. A size-$m$ rebuild uses $O(m\log^4(2+m))$ application work by
inserting its edges in parent-before-child order through the same source
interface, and allocates
$O(m\log^3(2+m))$ persistent hull state. The geometric sum of the doubled
capacities is $O(n)$, so these bounds remain $O(n\log^4(2+n))$ work and
$O(n\log^3(2+n))$ retained state in total. This alternative never reserves
capacity indexed by the original graph's unknown size.
\end{proof}

\begin{theorem}[Local exact obstacle and ACL solves on trees; \statusproved]
\label{thm:op3-local-top-tree}
On an arbitrary finite, simple, connected, unweighted tree with a supplied
seed vertex and local degree/adjacency queries, there exists a deterministic
algorithm returning the exact obstacle solution for
$\boldsymbol M=\boldsymbol D-\gamma\boldsymbol A$ and
$\boldsymbol b=\boldsymbol e_v-\lambda\boldsymbol d$.
If $U$ is its returned positive support, its total exact-real word work is
\begin{equation}
 O\!\left(1+\operatorname{cvol}(U)
     \log^4(2+\operatorname{cvol}(U))\right),
 \label{eq:op3-local-top-tree-work}
\end{equation}
and its space, even retaining every historical application version, is
$O(1+\operatorname{cvol}(U)\log^3(2+\operatorname{cvol}(U)))$ words.
For $0<\epsappr<1$ and $\lambda=\epsappr/2$, the vector
$\widehat{\boldsymbol\pi}=\bar\alpha\boldsymbol D\boldsymbol u$
has original source residual
$0\leq\boldsymbol e_v-\boldsymbol M\boldsymbol u
 \leq\lambda\boldsymbol d\leq\epsappr\boldsymbol d$.
Consequently its fully charged local ACL work is
$\widetilde O(1/\epsappr)$, with no polynomial dependence on $1/\alpha$.
\end{theorem}
\begin{proof}
First query $d_v$. If $\lambda d_v\geq1$, return zero, which is the exact
obstacle solution. Otherwise the singleton face has
$u_v=(1-\lambda d_v)/d_v>0$. Scan its row, query and cache neighbor degrees,
and build its heap of thresholds $\lambda d_j/\gamma$. The singleton's
minimum threshold gives the exact next gate or a quiet stop directly.
When the first child is admitted, create the first active edge and home
both endpoint rows there.

For every later positive face, retain the source-exposed top tree with its
complete root summary. The scalar Schur equation gives $u_v$ in constant
work. The root's minimum threshold either names an actual strict positive
original boundary gate or certifies all gates nonpositive, since one-port
normalizations preserve signs. A named candidate has exactly one active
parent: two active parents, together with the path between them in the
connected active set, would create a cycle in the original tree.

Admit the named candidate $j$. Remove it from its parent's local heap and
refresh that parent's one home-edge payload. Scan $j$'s original row once;
its sole previously active neighbor is its parent. Cache its other
neighbors and their degree replies, build its own heap and one home row,
and allocate its new active vertex and edge records. Invoke one source
\texttt{link}, then \texttt{expose} on $v$. Original degrees and loads in all
existing numerical records remain fixed. Changes to active incidence and
port sets are handled by the source's split/create/join discipline, while
\cref{lem:op3-source-interior-adapter} handles temporary source-interior
clusters. Only the heap minimum is homed on an edge; the other original
candidate thresholds remain in the local heap for later admissions.

Every admitted gate is evaluated at the exact restricted solution. A
positive Schur complement correction gives $u_j>0$ and increases all old
active coordinates, so the face stays positive and grows monotonically.
Summing the original residual over the entire graph gives
\[
 \boldsymbol1^{\mathsf T}
 (\boldsymbol e_v-\boldsymbol M\boldsymbol u)
 =1-\bar\alpha\boldsymbol d^{\mathsf T}\boldsymbol u.
\]
Every inactive residual is nonnegative, and each active residual equals
$\lambda d_i$. Therefore $\lambda\vol(U)<1$ at every positive checkpoint.
There are finitely many strict admissions. At a quiet checkpoint the active
equations and all boundary inequalities are exactly the obstacle KKT
conditions; other inactive rows have zero residual. Uniqueness gives the
exact obstacle solution and the displayed ACL certificate. This proof
identity does not require an algorithmic scan of the whole graph.

Let $V=\operatorname{cvol}(U)$ for the final nonempty support. Scanning every
admitted row once costs at most $V$. The number of discovered candidate
records and degree replies is at most $V$; dictionary accesses, heap
construction/deletion and all unsuccessful heap or final quiet checks cost
$O(V\log(2+V))$. Each admission uses one home-edge update, one link and one
expose, hence $O(\log(2+V))$ recomputed cluster records by the source theorem
and \cref{lem:op3-top-tree-payload}. Every callback costs at most
$O(\log^3(2+V))$ word work and allocates $O(\log^2(2+V))$ hull nodes.
Thus all comparisons, searches, factor updates, copies and historical
records satisfy \eqref{eq:op3-local-top-tree-work} and the claimed space
bound. Structural source state is linear in active size and already paid.
The optional capacity-doubling rebuilds in
\cref{lem:op3-top-tree-payload} satisfy these same total bounds.

Finally solve the root's scalar equation once and traverse only the current
cluster hierarchy with its saved affine recovery records. Each represented
edge or helper record is visited a constant number of times. Active indices
allow each vertex to be emitted once; comparison dictionaries would add only
an already allowed logarithm. Multiply by $\bar\alpha d_i$ for the sparse
PageRank output. Neither the past hierarchies nor old full solution vectors
are traversed. Since $V\leq2\vol(U)<2/\lambda$, the stated ACL bound follows.
\end{proof}

\paragraph{An exact RPPR consequence.}
In the nonzero regime $0<\rho<1/d_v$, taking $\lambda=\rho$ and returning
$\boldsymbol x=\bar\alpha\boldsymbol D^{1/2}\boldsymbol u$ gives the exact
shared RPPR minimizer on trees, by the mapping in \cref{sec:mapping}.
Its word work is $\widetilde O(1/\rho)$ and it has zero objective gap in this
exact-real model. This structural statement does not convert ACL accuracy
into the semantic PPR tolerance of OP1 and does not assert numerical stability.

\paragraph{Exact implementation evidence.}
\texttt{top\_tree\_callback\_adapter.py} implements the application layer.
Its final audit covers 650 canonical positive faces: every nonisomorphic
core tree through eight vertices, every seed, and two teleportation choices.
Each core vertex has one inactive private-star report. The original degrees
and tiny positive $\lambda$ give a legal positive core face with every report
quiet; the stars need not be materialized. All 167,384 legal binary cluster
joins under exposed and unexposed seed boundaries are checked, covering all
five join shapes. There are 26,130 independent Schur oracles, 168,294 valid
summary comparisons, 104,794 owned rows, 270,792 exact two-port queries,
78,030 one-port checks and 833,718 recovered coordinates. The audit also
checks 650 complete exposure transitions and 1,296 intermediate edge-payload
refreshes, then rechecks the old roots. Payload-only perturbations test the
algebraic interface; they are not represented as legal graph admissions.

The driver enumerates clusters, constructs supplied hierarchies and performs
dense reference solves. Its 240.431-second runtime does not measure the
source hierarchy's online bound. The asymptotic result imports that source
algorithm and proves every additional application cost above; it does not
infer complexity from finite timing. Audit counts, graph/parameter rules,
source and backend hashes are in \texttt{TOP\_TREE\_CALLBACK\_AUDIT.json}.
